\section{Conclusion}

This work developed a methodology for systematically characterizing the trade-offs between computational performance and reconstruction fidelity of Generalizable and Per-Scene NeRF architectures under sparse-view conditions. The experimental results demonstrate that Per-Scene models achieve substantially lower inference latency and VRAM consumption than generalizable models, making them well suited for deployment in resource-constrained environments. However, these benefits rely on dense multi-view supervision, restricting their usage in sparse-view acquisition scenarios, where generalizable architectures remain a practical alternative despite reduced reconstruction fidelity.

Although the proposed strategy provides a standardized and reproducible framework, some limitations should be acknowledged. The restricted hardware configuration of the experimental environment prevented the inclusion of architectures that rely on computationally intensive 3D cost volumes or strong geometric regularization, such as RegNeRF. Furthermore, the experimental validation was conducted on the LLFF dataset, and the adopted sparse-view sampling strategy exploits the sequential ordering of its image indices, which may limit its direct applicability to unordered image collections.

In this context, future work will focus on extending the proposed methodology to unbounded 360-degree datasets and more diverse acquisition scenarios, enabling a broader characterization of these trade-offs. This includes incorporating sparse-view 3D Gaussian Splatting variants \cite{seo2026improvingsparseview3dgsgeneralization,song2025d2gsdepthanddensityguidedgaussian} into the benchmarking protocol, to assess whether their explicit representation offers a more favorable fidelity--efficiency trade-off than the implicit architectures evaluated here. Also, diffusion priors will be investigated as regularization mechanisms during FT to improve reconstruction quality.